\documentclass[aspectratio=169,UTF8]{ctexbeamer}
\usetheme{Madrid}
\usecolortheme{default}
\definecolor{DeepBlue}{HTML}{2563EB}
\definecolor{DeepInk}{HTML}{172033}
\definecolor{DeepGreen}{HTML}{059669}
\setbeamercolor{structure}{fg=DeepBlue}
\setbeamercolor{title}{fg=white,bg=DeepBlue}
\setbeamercolor{frametitle}{fg=white,bg=DeepBlue}
\setbeamercolor{alerted text}{fg=DeepGreen}
\setbeamertemplate{navigation symbols}{}
\setbeamertemplate{footline}[frame number]
\usepackage{booktabs}
\usepackage{listings}
\usepackage{xcolor}
\usepackage{hyperref}
\hypersetup{colorlinks=true,urlcolor=DeepBlue}
\lstset{
  basicstyle=\ttfamily\scriptsize,
  keywordstyle=\color{DeepBlue},
  stringstyle=\color{DeepGreen},
  commentstyle=\color{gray},
  showstringspaces=false,
  breaklines=true,
  frame=single,
  columns=fullflexible
}
\newcommand{\api}[1]{\texttt{#1}}
\newcommand{\keypoint}[1]{\alert{#1}}
\AtBeginSection[]{
  \begin{frame}{本节路线}
    \tableofcontents[currentsection]
  \end{frame}
}

\title{第1课：认识 DeepSeek API 与首次调用}
\subtitle{从账号、鉴权到第一个可观测请求}
\author{DeepSeek API 全栈课程}
\date{文档快照：2026-11-05}
\begin{document}
\frame{\titlepage}
\begin{frame}{学习目标}
\begin{itemize}
  \item 理解 API、模型、Token、SDK 的关系
  \item 安全配置 API Key，并分别用 HTTP 与 OpenAI SDK 调用
  \item 选择 \api{deepseek-flash} / \api{deepseek-v4-pro}
  \item 查询模型、余额并读懂最小响应
\end{itemize}
\end{frame}
\begin{frame}{全课程能力阶梯}
\centering
鉴权与对话 $\rightarrow$ 流式与推理 $\rightarrow$ 结构化与工具\\[6pt]
$\rightarrow$ 视觉/文件/搜索 $\rightarrow$ 多协议兼容 $\rightarrow$ 生产 Agent
\end{frame}
\section{基本概念}
\begin{frame}{一次调用的组成}
\begin{block}{请求}
端点 + Bearer Key + 模型 + 输入消息 + 生成参数
\end{block}
\begin{block}{响应}
候选结果 + 停止原因 + Token 用量 + 请求标识
\end{block}
DeepSeek API 是\keypoint{无状态}服务：客户端负责保存需要延续的对话。
\end{frame}
\begin{frame}{当前模型选择}
\begin{tabular}{lll}
\toprule
模型 & 主要定位 & 图像输入\\
\midrule
\api{deepseek-flash} & 速度、成本、视觉 & 支持\\
\api{deepseek-v4-pro} & 高质量复杂任务 & 不支持\\
\bottomrule
\end{tabular}
\vspace{8pt}

两者支持 1M 上下文、思考模式、JSON、工具、Responses 与 Anthropic 格式。能力和价格会更新，应以官方模型页为准。
\end{frame}
\begin{frame}{地址与兼容层}
\begin{description}
  \item[OpenAI 格式] \api{https://api.deepseek.com}
  \item[Anthropic 格式] \api{https://api.deepseek.com/anthropic}
  \item[Beta 功能] \api{https://api.deepseek.com/beta}
\end{description}
“兼容”表示可复用 SDK 与请求形态，不代表所有上游字段都生效。第10、11课逐项核对差异。
\end{frame}
\section{鉴权与环境}
\begin{frame}[fragile]{Key 只放环境变量}
\begin{lstlisting}[language=bash]
# PowerShell：仅当前终端会话
$env:DEEPSEEK_API_KEY="sk-..."
$env:DEEPSEEK_MODEL="deepseek-flash"
\end{lstlisting}
\begin{alertblock}{禁止}
不要提交 Key、打印 Key、放进前端、日志或示例截图；泄漏后立即吊销并轮换。
\end{alertblock}
\end{frame}
\begin{frame}{HTTP 鉴权}
\api{Authorization: Bearer <API\_KEY>}\\[8pt]
JSON 请求还需 \api{Content-Type: application/json}。常见失败：
\begin{itemize}
  \item 401：Key 缺失、拼写错误、已失效
  \item 402：余额不足
  \item 400/422：JSON 或字段不合法
\end{itemize}
\end{frame}
\section{首次调用}
\begin{frame}[fragile]{OpenAI SDK 最小调用}
\begin{lstlisting}[language=Python]
client = OpenAI(
    api_key=os.environ["DEEPSEEK_API_KEY"],
    base_url="https://api.deepseek.com",
)
response = client.chat.completions.create(
    model="deepseek-flash",
    messages=[{"role": "user", "content": "你好"}],
    max_tokens=64,
)
print(response.choices[0].message.content)
\end{lstlisting}
\end{frame}
\begin{frame}{最小响应阅读法}
\begin{description}
  \item[\api{id}] 请求/结果标识，排障时记录
  \item[\api{choices[0].message}] 助手消息
  \item[\api{finish\_reason}] stop、length、tool\_calls 等
  \item[\api{usage}] 输入、输出、总 Token 与缓存明细
  \item[\api{system\_fingerprint}] 后端配置指纹
\end{description}
先检查停止原因，再消费内容，最后记录 usage。
\end{frame}
\begin{frame}{模型与余额端点}
\begin{itemize}
  \item \api{GET /models}：运行时发现模型及能力元数据
  \item \api{GET /user/balance}：读取余额可用状态与币种明细
  \item 余额是账号级；模型列表适合健康检查，不应每请求查询
\end{itemize}
\end{frame}
\section{成本与练习}
\begin{frame}{Token 与计费直觉}
总费用通常由\keypoint{未命中输入、缓存命中输入、输出}分别计价；Search/Fetch 按成功请求次数计费。\\[8pt]
实际 Token 以响应 \api{usage} 为准；中文字符、英文单词与 Token 并非固定一一对应。
\end{frame}
\begin{frame}{本课例程}
\begin{enumerate}
  \item \api{01\_first\_call.py}：SDK 最小调用与 usage
  \item \api{02\_raw\_http.py}：理解裸 HTTP 请求
  \item \api{03\_account\_info.py}：模型列表与余额
\end{enumerate}
练习：比较两个模型的延迟、输出与 Token；不要把单次结果当成基准测试。
\end{frame}
\begin{frame}{检查清单}
\begin{itemize}
  \item Key 未进入源码或版本库
  \item 显式设置超时与 \api{max\_tokens}
  \item 记录 id、model、finish\_reason、usage
  \item 不硬编码价格；上线前复核官方页面
\end{itemize}
\end{frame}
\begin{frame}{官方资料}
\small
\url{https://api-docs.deepseek.com/zh-cn/}\\
\url{https://api-docs.deepseek.com/zh-cn/quick_start/pricing}\\
\url{https://api-docs.deepseek.com/zh-cn/quick_start/token_usage}\\
\url{https://api-docs.deepseek.com/zh-cn/api/list-models}\\
\url{https://api-docs.deepseek.com/zh-cn/api/get-user-balance}
\end{frame}
\end{document}
